% diagrams/firstwords/sit.tex -- Lucía y Amy sentadas en clase, cada una
% en su silla, vistas de lado (la cabeza, de frente, como en todos los
% dibujos): la espalda recta, el muslo sobre el asiento, la pierna
% colgando y las manos en las rodillas.
\begin{tikzpicture}[dibujofw]
  \foreach \x/\quien in {0/Lucia, 3.6/Amy} {
    \begin{scope}[shift={(\x,0)}]
      % la silla, de lado: el respaldo, el asiento y las patas
      \filldraw[fill=white] (0,0) rectangle (0.22,4.5);
      \filldraw[fill=white, rounded corners=1mm] (-0.05,3.3) rectangle (0.3,4.55);
      \filldraw[fill=white] (1.8,0) rectangle (2.02,2.1);
      \filldraw[fill=white] (-0.05,2.1) rectangle (2.15,2.4);
      % la pierna, colgando, y el zapato
      \filldraw[fill=white] (1.95,0.75) rectangle (2.32,2.6);
      \filldraw[fill=white, rounded corners=1.5mm] (1.9,0.45) rectangle (2.8,0.8);
      % el cuello y el vestido, del hombro a la rodilla
      \filldraw[fill=white] (0.72,4.3) rectangle (1.06,4.95);
      \filldraw[fill=white, rounded corners=2.5mm] (0.4,4.6) -- (1.35,4.6) -- (1.45,3.35)
        -- (2.4,3.05) -- (2.4,2.4) -- (0.3,2.4) -- cycle;
      % el brazo, con la mano en la rodilla
      \filldraw[fill=white, rounded corners=2mm] (0.7,4.45) -- (1.75,3.3) -- (1.55,3.05)
        -- (0.6,3.95) -- cycle;
      \filldraw[fill=white] (1.85,3.2) circle (0.22);
      % la cabeza (diagrams/firstwords/kit.tex)
      \begin{scope}[shift={(0.89,5.75)}, scale=0.85] \csname cabeza\quien\endcsname \end{scope}
    \end{scope}
  }
  % el suelo
  \draw (-0.5,0) -- (6.9,0);
\end{tikzpicture}
